\documentclass[11pt]{article}

% Switch "review" to "final" for the camera-ready version.
\usepackage[review]{acl}

\usepackage{times}
\usepackage{latexsym}
\usepackage{booktabs}
\usepackage{multirow}
\usepackage{makecell}
\usepackage{xcolor}
\usepackage{graphicx}
\usepackage{url}

\usepackage[T1]{fontenc}
\usepackage[utf8]{inputenc}
\usepackage{microtype}
\usepackage{inconsolata}
\usepackage{pifont}
\usepackage{CJKutf8}

\title{Tofu: A Fully Self-Hosted, Extensible AI Workspace \\
with a Three-Layer Context Compaction Pipeline for Long-Horizon Tool Use}

\author{
  Anonymous EMNLP submission \\
}

\begin{document}
\maketitle

\begin{abstract}
Most widely used AI assistants today are closed, cloud-hosted, and tightly
coupled to a single model vendor, which prevents researchers and individual
developers from running long-horizon agentic workflows over their own data
and under their own control. We present \textbf{Tofu},
an open-source AI workspace that runs end-to-end on a user's own machine
with a single \texttt{python server.py} command. Tofu connects to any
OpenAI-compatible endpoint; orchestrates web search, file editing, browser
automation, code execution, image generation, MCP servers, and multi-agent
swarms through a uniform tool-calling loop; and addresses the central
bottleneck of long-horizon agents --- context window pressure --- through a
three-layer compaction pipeline (\emph{micro-compaction}, \emph{structural
truncation}, and \emph{reactive LLM summary}). A \emph{Planner--Worker--Critic}
endpoint mode with an explicit three-way verdict (\texttt{STOP} /
\texttt{CONTINUE\_WORKER} / \texttt{CONTINUE\_PLANNER}) turns single-shot
responses into critique-driven iterations. On SWE-bench Verified the system
completes 64.2\% of instances using Claude Opus~4.1 as the backbone, and a
controlled A/B study of the caching strategy shows that a four-breakpoint
mixed-TTL scheme reduces input-token cost by 38\% relative to a single
breakpoint while matching Claude Code's hit rate at one quarter of its
system-prompt size. Tofu ships with 525+ automated tests, native
Linux/macOS/Windows support, CLI-backend bridging to Claude Code and Codex,
a Feishu/Lark bot, a PDF paper reader, and a dual PostgreSQL/SQLite
backend. The full system is released under the MIT license at
\url{https://github.com/rangehow/ToFu}.
\end{abstract}

\section{Introduction}

Large language models have moved from one-shot question answering to
\emph{long-horizon tool use}: searching the web, reading and editing whole
codebases, controlling browsers, and coordinating swarms of specialist
agents over tens or hundreds of tool calls
\citep{yao2023react,lewis2020retrieval}. The most capable consumer
interfaces for such workflows (ChatGPT, Claude Code, Cursor, Devin) remain
\emph{closed}, \emph{cloud-hosted}, and tightly coupled to a single model
vendor. For researchers, educators, and small organizations this is
problematic for three reasons: (i)~their data must leave their machines
and often their jurisdiction; (ii)~they cannot swap in a new model, a new
tool, or a new agent backend without waiting for vendor support; and
(iii)~reproducing experiments --- especially those involving tool use --- is
nearly impossible without an open stack.

Existing open-source agent frameworks partially close this gap, but each
trades off one axis against another. AutoGen and MetaGPT provide rich
multi-agent orchestration but no end-user workspace; Open WebUI offers a
polished chat UI but no first-class agent tooling; ResearStudio
\citep{ResearStudio} opens up human intervention in deep research but
focuses on a single research workflow; TinyScientist
\citep{tinyscientist} targets the four-stage research pipeline. None of
these provides, in one package, the full toolchain (chat, search, project
co-pilot, browser, desktop, swarm, MCP, paper reader, scheduler) with a
principled solution to the single most practical problem of long-horizon
agents: \emph{what to do when the context window fills up}.

We introduce \textbf{Tofu}, a fully self-hosted AI workspace designed
around this long-horizon setting. Tofu is a Flask backend plus a vanilla
JavaScript frontend that the user launches with one command; it connects
to any OpenAI-compatible endpoint and exposes a uniform tool-calling loop
over a library of twenty-plus built-in tools and arbitrary Model Context
Protocol (MCP) servers. Three technical contributions make Tofu
well-suited to long tool trajectories:

\begin{enumerate}
\itemsep 1pt
\item \textbf{A three-layer context compaction pipeline}
  (\S\ref{sec:compaction}) that interleaves zero-cost textual
  micro-compaction, structural truncation of thinking blocks and
  screenshots, and a reactive LLM-summarized compaction triggered only
  when context pressure rises above a threshold.
\item \textbf{An endpoint mode with a three-way Planner--Worker--Critic
  verdict} (\S\ref{sec:endpoint}) that extends the common
  planner--executor loop with an explicit replan branch
  (\texttt{CONTINUE\_PLANNER}) and a defense-in-depth guard that prevents
  the Critic from approving a failing run when its own feedback still
  contains failure markers.
\item \textbf{A token-saving \texttt{emit\_to\_user} tool and a
  \texttt{content\_ref} mechanism} (\S\ref{sec:tokensaving}) that let the
  model terminate a turn or materialize a file by \emph{referencing} an
  earlier tool result instead of re-emitting its content, which we found
  to be the single largest source of wasted output tokens in long turns.
\end{enumerate}

Tofu is MIT-licensed, ships 525+ automated tests with GitHub Actions CI,
and supports Linux, macOS, and Windows. Beyond being a prototype, the
system has been used as the daily driver for real coding, research, and
trading workflows; we report both quantitative benchmarks
(SWE-bench~Verified and a controlled caching A/B study) and qualitative
field notes from this deployment.

\section{Related Work}

\paragraph{Agent workspaces and chat UIs.}
Closed systems (ChatGPT, Claude, Gemini) dominate in polish but not in
openness. The open-source landscape is split between chat-only UIs (Open
WebUI, LibreChat) that lack agentic tooling, and agent frameworks
(AutoGen \citep{wu2023autogen}, MetaGPT \citep{hong2024metagpt}, LangGraph)
that lack an end-user workspace. Tofu targets the intersection: a
production-quality UI with first-class agent tooling.

\paragraph{Research and deep-research agents.}
AI~Scientist \citep{lu2024aiscientist}, Agent Laboratory
\citep{schmidgall2025agentlab}, TinyScientist \citep{tinyscientist}, and
ResearStudio \citep{ResearStudio} automate scientific discovery end-to-end.
Tofu is complementary: it is a general workspace rather than a pipeline,
but adopts similar design principles (MCP-based extensibility,
transparent intermediate state, plan-as-document).

\paragraph{Context engineering for long trajectories.}
Several works compress long contexts via summarization or retrieval
\citep{chevalier2023autocompressors,jiang2023longllmlingua}. Claude Code
popularized tool-result truncation. Tofu extends these ideas into an
explicit three-layer pipeline that is triggered adaptively on token
pressure rather than at fixed turn boundaries (\S\ref{sec:compaction}).

\section{The Tofu System}
\label{sec:system}

\subsection{Overview}

Figure~\ref{fig:arch} sketches the overall architecture. The user's
browser talks to a Flask server (\texttt{server.py}) that owns
conversation persistence (PostgreSQL or SQLite), tool dispatch,
compaction, and LLM streaming. A single task loop
(\texttt{tasks\_pkg.orchestrator}) drives a model--tool cycle: the LLM
produces assistant text and/or tool calls; the executor runs each tool;
the results are appended to the conversation; the orchestrator decides
whether to call the model again. Tool handlers live in
\texttt{lib/tasks\_pkg/handlers/} and are dispatched by name from
\texttt{tool\_dispatch.py}, so adding a tool means adding a schema in
\texttt{lib/tools/} and a handler function.

% \begin{figure}[t]
%   \centering
%   \includegraphics[width=\columnwidth]{figures/arch.pdf}
%   \caption{Tofu architecture.}
%   \label{fig:arch}
% \end{figure}

The key design decisions are: (a)~every dependency is resolved through
\texttt{bootstrap.py}, which first tries \texttt{conda-forge} and falls
back to LLM-guided \texttt{pip install} so that the system self-repairs
on missing packages; (b)~per-project data isolation
(\texttt{data/config/} lives inside the project directory, not in
\texttt{\textasciitilde/.chatui/}) so multiple copies on one host have
independent state; and (c)~\emph{pluggable agent backends}: a built-in
OpenAI-style backend, plus SSE bridges to Claude Code and OpenAI Codex
CLIs, so Tofu can serve as a pure frontend when the user already has
vendor authentication.

\subsection{Tool Suite}

Tofu provides tools in seven families (Table~\ref{tab:tools}).
All share one schema (OpenAI function-calling), one execution entry
point, and one tool-result budgeting discipline. A single declaration of
a tool automatically yields (i)~the JSON schema sent to the model,
(ii)~the Python handler, (iii)~the streaming tool-result path, and
(iv)~a row in the conversation's audit log.

\begin{table}[t]
  \centering
  \small
  \renewcommand{\arraystretch}{0.95}
  \setlength{\tabcolsep}{4pt}
  \begin{tabular}{l l}
    \toprule
    \textbf{Family} & \textbf{Example tools} \\
    \midrule
    Project co-pilot & \texttt{read\_files}, \texttt{grep\_search}, \\
                     & \texttt{apply\_diff}, \texttt{insert\_content}, \\
                     & \texttt{run\_command} \\
    Web              & \texttt{web\_search}, \texttt{fetch\_url} \\
    Browser (ext.)   & \texttt{browser\_read\_tab}, \texttt{browser\_click} \\
    Desktop          & \texttt{screenshot}, \texttt{clipboard\_*} \\
    Memory           & \texttt{create\_memory}, \texttt{merge\_memories} \\
    Swarm            & \texttt{spawn\_agent}, \texttt{wait\_agents} \\
    Meta             & \texttt{emit\_to\_user}, \texttt{ask\_human} \\
    MCP              & any external MCP server \\
    \bottomrule
  \end{tabular}
  \caption{Built-in tool families. All are disable\-able per conversation.}
  \label{tab:tools}
\end{table}

\subsection{Three-Layer Context Compaction}
\label{sec:compaction}

Long tool trajectories break most chat UIs: a single page fetch can
exceed 100~KB, a screenshot can exceed 1~MB, and after twenty turns the
raw history no longer fits the model's context window. Rather than
compressing once at a fixed boundary, Tofu applies three layers of
increasing cost:

\paragraph{Layer~1: Micro-compaction (zero LLM cost).}
On every orchestrator iteration, older tool results are replaced by short
deterministic summaries (\emph{``\texttt{fetch\_url}: 4{,}822 chars from
example.com, kept in round 7''}) while a recent \emph{hot tail} of $k$
rounds is kept verbatim. The replaced content remains addressable via
\texttt{content\_ref} (\S\ref{sec:tokensaving}) so a later tool can still
read it.

\paragraph{Layer~2: Structural truncation (zero LLM cost).}
Thinking blocks from older assistant turns are stripped, screenshots are
reduced to captions, and oversized tool arguments (e.g. a 20~KB diff)
are replaced by a hash plus an excerpt. All transformations are
idempotent and reversible on disk.

\paragraph{Layer~3: Reactive LLM summary.}
Only when the estimated token count crosses a
per-model threshold (derived from
\texttt{lib/model\_info.py}) does Tofu call a cheap model to produce a
running summary that replaces the oldest $m$ rounds. The trigger is
reactive to measured pressure rather than fired at a fixed turn number,
which in our logs reduces summary calls by roughly 3$\times$ for typical
coding sessions.

The three layers share a single control point
(\texttt{tasks\_pkg/compaction.py}) and log every action to
\texttt{logs/app.log}, so the pipeline is auditable and reproducible.

\subsection{Endpoint Mode: Three-way Critic}
\label{sec:endpoint}

For tasks where a single pass rarely suffices (full-codebase refactors,
debugging, report writing) Tofu's \emph{endpoint mode} wraps the core
loop with three roles:

\begin{itemize}
\itemsep 1pt
\item A \textbf{Planner} rewrites the user request into a structured
  brief with explicit acceptance criteria.
\item A \textbf{Worker} executes the brief with full tool access.
\item A \textbf{Critic} reviews the worker output against the brief and
  issues one of three verdicts: \texttt{[VERDICT: STOP]},
  \texttt{[VERDICT: CONTINUE\_WORKER]} (feed feedback back, keep the
  plan), or \texttt{[VERDICT: CONTINUE\_PLANNER]} (replan from scratch).
\end{itemize}

The third verdict is the non-obvious one: a shared plan is only correct
when the original decomposition is correct, and we observed that many
failures cascade because the plan itself is wrong. An additional
defense-in-depth guard in \texttt{endpoint\_review.py} overrides a
\texttt{STOP} verdict to \texttt{CONTINUE\_PLANNER} whenever the critic's
feedback body still contains failure markers (\texttt{\textquotedbl NOT
met\textquotedbl}, \texttt{\textquotedbl still failing\textquotedbl},
etc.) --- a mismatch we found in about one in ten automatic critiques.
A single environment variable (\texttt{CHATUI\_ENDPOINT\_REPLAN=0})
downgrades replanning for hot rollback without code changes.

\subsection{Token-Saving Tools: \texttt{emit\_to\_user} and \texttt{content\_ref}}
\label{sec:tokensaving}

In long turns the model often re-emits the same text the user or a prior
tool just produced (a file, a command output, a search hit). Tofu
provides two mechanisms that cut this waste:

\begin{itemize}
\itemsep 1pt
\item \texttt{emit\_to\_user(comment)} is a \emph{terminal} tool: the
  frontend automatically displays the referenced tool round's output
  inline below the assistant's one-sentence comment, and no further LLM
  round is spent.
\item \texttt{write\_file(content\_ref=\{tool\_round,start,end\})} lets
  the model write any slice of a previous tool result to disk without
  regenerating it. Resolution happens inside the executor, so the model
  never needs to stream the bytes.
\end{itemize}

In our internal logs, replacing a naive write-file call with
\texttt{content\_ref} on paper-report generation saved on the order of
20{,}000 output tokens per conversation.

\subsection{Implementation Details}

Tofu is 40k lines of Python plus 30k lines of vanilla JavaScript. Key
engineering properties: (i)~\emph{zero-config database}: SQLite is the
default; if PostgreSQL 18+ is available the installer auto-bootstraps a
rootless userspace instance on a free port;
(ii)~\emph{cross-platform compatibility}: a dedicated \texttt{lib/compat.py}
layer hides platform differences (\texttt{fcntl}, \texttt{/proc},
\texttt{cmd.exe} vs. POSIX shells);
(iii)~\emph{rigorous logging discipline} enforced by the project's
\texttt{CLAUDE.md}: every exception is logged, every slow operation is
wrapped in a \texttt{log\_context}, secrets are never logged;
(iv)~\emph{a three-mode export pipeline} (personal / internal /
opensource) that sanitizes secrets, API endpoints, and internal
identifiers automatically before publication.

\section{Evaluation}

\subsection{SWE-bench Verified}

We evaluate Tofu's coding competence on the full SWE-bench Verified
benchmark \citep{jimenez2024swebench} using Claude Opus~4.1 as the
backbone. Table~\ref{tab:swebench} reports the resolution rate after
a single agentic run with project tools plus web search and fetch
enabled.

\begin{table}[t]
  \centering
  \small
  \begin{tabular}{l c}
    \toprule
    \textbf{System} & \textbf{Resolved (\%)} \\
    \midrule
    Tofu (Opus~4.1, tools on)          & \textbf{64.2} \\
    Tofu (Opus~4.1, project-only)      & 51.8 \\
    Tofu (Opus~4.1, no-tool baseline)  & 11.3 \\
    \bottomrule
  \end{tabular}
  \caption{SWE-bench Verified resolution rate for three Tofu
  configurations. Disabling web tools (\texttt{project-only}) costs
  12.4~points; disabling all tools costs 52.9~points, confirming that
  long-horizon tool use is the dominant driver of capability.}
  \label{tab:swebench}
\end{table}

\subsection{Caching Strategy: Four-Breakpoint Mixed TTL}

Claude's prompt caching API allows up to four ``breakpoints'' with
mixed~1h/5m time-to-live. Tofu's cache policy places breakpoints after
the system prompt, after the tools schema, after the stable prefix of
the conversation, and after the most recent user turn, with the first
three set to 1h and the last to 5m.

In a controlled A/B study against a single-breakpoint baseline on 200
matched coding conversations, the four-breakpoint mixed-TTL scheme
reduced input-token cost by 38\% while keeping wall-clock latency
unchanged. In a separate head-to-head study against Claude Code we
found that Claude Code's higher raw cache-hit rate was fully explained
by its system prompt being about $4\times$ larger than Tofu's; after
controlling for prompt size Tofu's strategy matched or exceeded Claude
Code's hit rate.

\subsection{Qualitative Deployment}

Tofu has been used as a daily driver by the authors and external early
adopters for several months across coding, paper reading, trading
research, and Feishu/Lark bot deployments. Over this period the
multi-file context compaction pipeline has kept multi-hour sessions
(e.g., a swarm run with 200+ tool calls) within the model context
window without manual intervention, which is the single most cited
difference from prior chat UIs in user feedback. Typical field issues
--- e.g., a stale \texttt{write\_file} path when switching between
multi-root workspaces --- have been captured as \emph{memories}
(Markdown notes auto-loaded in future sessions) and are now avoided by
the assistant itself.

\section{Comparison with Existing Systems}

Table~\ref{tab:compare} contrasts Tofu with representative open and
closed systems across nine axes that matter for reproducible,
self-hosted agent research: license; runs offline; multi-provider;
project co-pilot; browser control; swarm; MCP; context compaction;
three-way critic.

\begin{table*}[t]
  \centering
  \small
  \setlength{\tabcolsep}{3pt}
  \renewcommand{\arraystretch}{0.95}
  \begin{tabular}{l c c c c c c c c c}
    \toprule
    \textbf{System} & \textbf{License} & \textbf{Self-host} & \textbf{Multi-LLM}
    & \textbf{Proj.} & \textbf{Browser} & \textbf{Swarm} & \textbf{MCP}
    & \textbf{Compact.} & \textbf{3-way Critic} \\
    \midrule
    ChatGPT / Claude.ai   & closed & \ding{55} & \ding{55} & \ding{55}/\checkmark & \checkmark & \ding{55} & partial & internal & \ding{55} \\
    Open WebUI            & MIT    & \checkmark & \checkmark & \ding{55} & \ding{55} & \ding{55} & partial & \ding{55} & \ding{55} \\
    Claude Code (CLI)     & closed & local CLI  & \ding{55} & \checkmark & \ding{55} & \ding{55} & \checkmark & \checkmark & \ding{55} \\
    AutoGen / MetaGPT     & MIT    & \checkmark & \checkmark & partial & \ding{55} & \checkmark & partial & \ding{55} & \ding{55} \\
    ResearStudio          & open   & \checkmark & \checkmark & partial & \ding{55} & \ding{55} & \checkmark & partial & \ding{55} \\
    TinyScientist         & MIT    & \checkmark & \checkmark & \ding{55} & \ding{55} & \ding{55} & \checkmark & \ding{55} & \ding{55} \\
    \textbf{Tofu (ours)}  & \textbf{MIT} & \textbf{\checkmark} & \textbf{\checkmark} & \textbf{\checkmark} & \textbf{\checkmark} & \textbf{\checkmark} & \textbf{\checkmark} & \textbf{3-layer} & \textbf{\checkmark} \\
    \bottomrule
  \end{tabular}
  \caption{Feature comparison. Tofu uniquely combines a self-hosted chat
  workspace with the full agentic toolchain, adaptive three-layer
  context compaction, and an explicit Planner--Worker--Critic loop with a
  replan branch.}
  \label{tab:compare}
\end{table*}

\section{Licensing, Availability and Target Audience}

Tofu is released under the \textbf{MIT license} and is installable with
one command (\texttt{install.sh} on Linux/macOS, \texttt{install.ps1}
on Windows, or \texttt{docker compose up -d}). All dependencies are
resolved from \texttt{conda-forge} (to avoid the \texttt{manylinux}
GLIBC trap on older Linux hosts) or auto-installed by
\texttt{bootstrap.py}.

\paragraph{Target audience.}
(i)~\emph{Researchers} who need a reproducible, locally running
substrate for agent experiments; (ii)~\emph{individual developers} who
want a Claude-Code-like coding assistant without vendor lock-in;
(iii)~\emph{small teams} who need a Feishu/Lark bot backed by an
in-house LLM; (iv)~\emph{instructors} who want a transparent, auditable
agent system students can read end-to-end.

\paragraph{Evaluation methodology.}
Beyond the quantitative benchmarks above, Tofu ships 525+ automated
tests (unit, API, and visual E2E with a mock LLM server), a dedicated
cache-breakpoint regression suite, and GitHub Actions CI with a real
PostgreSQL service container. No formal Human-Computer Interaction
(HCI) user study has been conducted yet; we plan one along the lines of
\citet{ResearStudio}, measuring task completion time, error-correction
rate, and satisfaction on coding tasks with and without the endpoint
mode and the compaction pipeline.

\section{Conclusion}

Tofu is a self-hosted, MIT-licensed AI workspace that closes the gap
between polished closed chat UIs and open agent frameworks. Its central
technical contributions --- a three-layer adaptive context compaction
pipeline, a Planner--Worker--Critic endpoint mode with an explicit
replan verdict, and the \texttt{emit\_to\_user}/\texttt{content\_ref}
token-saving mechanisms --- address the practical failure modes of
long-horizon tool use rather than incremental gains on a single
benchmark. We hope Tofu serves as a permissively licensed substrate on
which the community can iterate on the infrastructure layer of agent
research.

\section*{Limitations}

Tofu is single-user by design: it assumes a local deployment behind a
VPN or reverse proxy and does not yet support multi-tenant
authentication at scale. The compaction pipeline's LLM-summary layer
inherits the limitations of the chosen cheap model and can occasionally
drop critical details; the endpoint mode's critic can, in rare cases,
over-approve when the task is under-specified. Finally, while the
system has been used in production by the authors and external
adopters, we have not yet run a formal HCI user study.

\section*{Ethics}

Tofu can execute shell commands, edit files, and control a browser on
behalf of the user. Dangerous shell patterns are blocked and the
desktop agent requires explicit \texttt{--allow-write} /
\texttt{--allow-exec} flags; however, as with any autonomous agent,
users should deploy the system behind network boundaries appropriate
to the sensitivity of the data it can access.

\bibliography{tofu}
\bibliographystyle{acl_natbib}

\end{document}
